\section{Abstract}

Cryptoasset markets combine centralized exchanges (CEX) with automated market makers (AMM), whose price depends on reserves and on the trades executed against them. In live markets, arbitrage is constrained by costs, capital and liquidity; this study represents its capacity with a specified per-tick limit and does not model fees or binding capital constraints. It examines how a constant-product AMM tracks an exogenous CEX price after a shock, in a micro-world with noise, trend, arbitrage and liquidity-provision agents, and evaluates a two-variable discrete dynamical map as a reduced description. Three hypotheses are tested under a confirmatory protocol fixed in advance: confirmation seeds and cells disjoint from the exploration ones, inference paired by seed, and Holm correction.

With 300 confirmation seeds and 13500 runs, the specified active-arbitrage rule lowers the excess-error area from 2.36 to 0.10 within the tested simulator (H1; mean difference of −2.26; Holm p of 0.00015), and capacity grades recovery: 56.8 percent with low capacity and 100 percent with high capacity. Adaptive liquidity withdrawal raises the maximum response by 0.011 on average (H2), but only at the shallowest depth; at greater depths the effect is null. The reduced map beats persistence (H3), though not a scalar autoregression, and its out-of-sample regime classification is weak. A supplementary 40,500-path sensitivity found that 31.2 percent of recovered adaptive-LP paths later left the band; the shallow-depth amplification vanished when order sizes scaled with reserves and fell from 0.0346 to 0.0147 log units when panic was disabled. These conditional synthetic results suggest candidate measurements and questions for future real-data event-replay and stress-test studies; they do not validate monitoring procedures or operating thresholds. A separate six-minute data-readiness pilot contains three public DEX swaps; it is descriptive only and does not test the hypotheses. The confirmatory findings lack external validation, and the H1 precision target was not met.

\textbf{Keywords:} agent-based simulation; arbitrage; liquidity; price tracking; decentralized exchanges; discrete dynamical systems
